\subsection{二重逆极限}

设 I, L 是两个预序集，且 \(I \times L\) 为它们的乘积（§1，第4号）。考虑一个集合逆系统 \((\mathbf{E}_{\alpha}^{\lambda}, f_{\alpha\beta}^{\lambda\mu})\) （相对于指标集 \(I \times L\) ）。

我们有

\[
f _ {\alpha \gamma} ^ {\lambda \nu} = f _ {\alpha \beta} ^ {\lambda \mu} \circ f _ {\beta \gamma} ^ {\mu \nu} \quad \text { 当 } \alpha \leqslant \beta \leqslant \gamma \text { 且 } \lambda \leqslant \mu \leqslant \nu . \tag {11}
\]

让 \(E\) 或 \(\lim E_{\alpha}^{\lambda}\) 表示这个逆系统的逆极限。

\[
\overleftarrow {\alpha , \lambda}
\]

对于每一个 \(\stackrel{\alpha,}{\lambda} \in \mathbf{L}\) ，令 \(g_{\alpha \beta}^{\lambda} = f_{\alpha \beta}^{\lambda}: \mathbf{E}_{\beta}^{\lambda} \to \mathbf{E}_{\alpha}^{\lambda}\) 。由此从(11)可得

\[
g _ {\alpha \gamma} ^ {\lambda} = g _ {\alpha \beta} ^ {\lambda} \circ g _ {\beta \gamma} ^ {\lambda} \quad \text { 当 } \alpha \leqslant \beta \leqslant \gamma ,\tag{12}
\]

使得 \((\mathbf{E}_{\alpha}^{\lambda}, g_{\alpha \beta}^{\lambda})\) 是关于I的集合逆系统。设 \(\mathbf{F}^{\lambda} = \varprojlim_{\alpha} \mathbf{E}_{\alpha}^{\lambda}\) 记其逆极限。对于固定的 \(\lambda \leqslant \mu\) 在 L 中，由 (11) 可知 \(h_{\alpha}^{\lambda \mu} = f_{\alpha \alpha}^{\lambda \mu}: \mathrm{E}_{\alpha}^{\mu} \to \mathrm{E}_{\alpha}^{\lambda}\) 构成一个映射逆系统，其逆极限我们记为 \(h^{\lambda \mu} = \varprojlim_{\alpha} h_{\alpha}^{\lambda \mu}: \mathrm{F}^{\mu} \to \mathrm{F}^{\lambda}\) 。如果 \(\lambda \leqslant \mu \leqslant \nu\) 在 L 中，我们有

\[
h ^ {\lambda \nu} = h ^ {\lambda \mu} \circ h ^ {\mu \nu}\tag{13}
\]

（第 2 号，命题 2，推论 2）；因此 \((\mathbf{F}^{\lambda}, h^{\lambda \mu})\) 是关于 L 的集合逆系统。设其逆极限为 \(\mathbf{F} = \lim_{\lambda \in \mathbf{L}} \mathbf{F}^{\lambda}\) 。我们将定义一个典范双射 \(\mathbf{F} \to \mathbf{E}\) 。为此，注意到按照定义， \(\mathbf{F}\) 是 \(\prod_{\lambda \in \mathbf{L}} \mathbf{F}^{\lambda}\) 的子集，而 \(\mathbf{F}^{\lambda}\) 是 \(\prod_{\alpha \in \mathbf{I}} \mathbf{E}_{\alpha}^{\lambda}\) 的子集；因此可把 \(\mathbf{F}\) 典范地等同于 \(\prod_{(\alpha, \lambda) \in \mathbf{I} \times \mathbf{L}} \mathbf{E}_{\alpha}^{\lambda} = \mathbf{G}\) 的一个子集（第二章，§5，第5号，命题7）。对每个 \(z \in \mathbf{G}\) ，令 \(\operatorname{pr}^{\lambda}(z)\) 表示元素 \((\operatorname{pr}_{\alpha}^{\lambda}(z))_{\alpha \in \mathbf{I}}\) （属于 \(\prod_{z \in \mathbf{I}} \mathbf{E}_{\alpha}^{\lambda}\) ）。那么 \(z \in \mathbf{F}\) 当且仅当

\[
\operatorname{pr} ^ {\lambda} (z) = h ^ {\lambda \mu} (\operatorname{pr} ^ {\mu} (z)) \quad \text { 当 } \lambda \leqslant \mu \text { 中 } L\tag{14}
\]

并且 \(\mathbf{pr}^{\lambda}(z) \in \mathbf{F}^{\lambda}\) 对于所有 \(\lambda \in \mathbf{L}\) ；也就是说，每当 \(\alpha \leqslant \beta\) 在 I 中我们有

\[
\operatorname{pr} _ {\alpha} ^ {\lambda} (z) = f _ {\alpha \beta} ^ {\lambda \lambda} (\operatorname{pr} _ {\beta} ^ {\lambda} (z)).\tag{15}
\]

但 \(h^{\lambda \mu}(\mathrm{pr}^{\mu}(z)) = (f_{\alpha \alpha}^{\mu \lambda}(\mathrm{pr}_{\alpha}^{\mu}(z))_{\alpha \in \mathbf{I}}\) ；因此，由（14）和（15）可知，如果 \(\alpha \leqslant \beta\) 并且 \(\lambda \leqslant \mu\) ，我们有

\[
\operatorname{pr} _ {\alpha} ^ {\lambda} (z) = f _ {\alpha \alpha} ^ {\lambda \mu} (f _ {\alpha \beta} ^ {\mu \mu} (\operatorname{pr} _ {\beta} ^ {\mu} (z))) = f _ {\alpha \beta} ^ {\lambda \mu} (\operatorname{pr} _ {\beta} ^ {\mu} (z)),
\]

这意味着 \(z \in \mathbf{E}\) 。逆命题是显然的，因此我们已证明完毕。

命题4。如果 \((\mathbf{E}_{\alpha}^{\lambda}, f_{\alpha \beta}^{\lambda \mu})\) 是相对于预序集的乘积 \(\mathbf{I} \times \mathbf{L}\) 的集合逆系统，那么（在典范双射的意义下）我们有

\[
\varprojlim_ {\alpha , \lambda} \mathrm{E} _ {\alpha} ^ {\lambda} = \varprojlim_ {\lambda} (\varprojlim_ {\alpha} \mathrm{E} _ {\alpha} ^ {\lambda}).\tag{16}
\]

推论1. 设 \((\mathbf{E}_{\alpha}^{\prime \lambda}, f_{\alpha \beta}^{\prime \mu \lambda})\) 是相对于 \(\mathbf{I} \times \mathbf{L}\) 的另一个集合逆系统，并且对于每个 \((\alpha, \lambda) \in \mathbf{I} \times \mathbf{L}\) 让 \(u_{\alpha}^{\lambda}\) 为从 \(\mathbf{E}_{\alpha}^{\lambda}\) 到 \(\mathbf{E}_{\alpha}^{\prime \lambda}\) 中的一个映射，使得各 \(u_{\alpha}^{\lambda}\) 构成一个映射逆系统。那么

\[
\varprojlim_ {\alpha , \lambda} u _ {\alpha} ^ {\lambda} = \varprojlim_ {\lambda} (\varprojlim_ {\alpha} u _ {\alpha} ^ {\lambda}).\tag{17}
\]

验证过程与命题4的验证类似。

推论2。设 \((\mathrm{E}_{\alpha}^{\lambda}, f_{\alpha \beta}^{\lambda})_{\lambda \in \mathbf{L}}\) 是相对于 I 的一族集合逆系统。如果 \(\prod_{\lambda \in \mathbf{L}} f_{\alpha \beta}^{\lambda}\) 表示映射族 \((f_{\alpha \beta}^{\lambda})_{\lambda \in \mathbf{L}}\) 到乘积的扩展（第二章，第5节，第7号，定义2），那么 \(\left(\prod_{\lambda \in \mathbf{L}} \mathrm{E}_{\alpha}^{\lambda}, \prod_{\lambda \in \mathbf{L}} f_{\alpha \beta}^{\lambda}\right)\) 是关于指标集 I 的一个集合逆系统，并且（在典范双射的意义下）我们有

\[
\varprojlim_ {\alpha} \prod_ {\lambda \in \mathbf {L}} \mathrm{E} _ {\alpha} ^ {\lambda} = \prod_ {\lambda \in \mathbf {L}} (\varprojlim_ {\alpha} \mathrm{E} _ {\alpha} ^ {\lambda}).\tag{18}
\]

考虑二重逆系统 \((\mathbf{E}_{\alpha}^{\lambda}, g_{\alpha\beta}^{\lambda\mu})\) （相对于 \(I \times L\) ），其中L上的序关系是相等关系（第1号，例1），并应用命题4。

